% !TeX program = pdflatex
% Original architecture-engineering guidance; no licensed source text embedded.
\documentclass[11pt,letterpaper]{article}
\usepackage[T1]{fontenc}
\usepackage{lmodern,microtype}
\usepackage[margin=0.85in,headheight=16pt]{geometry}
\usepackage[table]{xcolor}
\usepackage{booktabs,tabularx,longtable,array,ragged2e}
\usepackage{tikz}
\usetikzlibrary{positioning,arrows.meta}
\usepackage{enumitem,needspace,fancyhdr,hyperref}
\definecolor{navy}{HTML}{17364D}
\definecolor{teal}{HTML}{087D83}
\definecolor{muted}{HTML}{516371}
\hypersetup{colorlinks=true,linkcolor=teal,urlcolor=teal}
\setlength{\parindent}{0pt}
\setlength{\parskip}{5pt}
\setlength{\emergencystretch}{3em}
\raggedbottom
\setlist{leftmargin=1.5em,itemsep=2pt,topsep=3pt}
\renewcommand{\arraystretch}{1.2}
\newcolumntype{Y}{>{\RaggedRight\arraybackslash}X}
\pagestyle{fancy}\fancyhf{}
\fancyhead[L]{\small\sffamily\color{navy} ARCHITECTURE ENGINEERING}
\fancyhead[R]{\small\sffamily\color{muted} SUITE / 1.0}
\fancyfoot[L]{\small\sffamily\color{muted} V\&B / ISO 42010:2022}
\fancyfoot[R]{\small\thepage}
\newcommand{\docsection}[1]{\Needspace{5\baselineskip}\section{#1}}
\newcommand{\docsub}[1]{\Needspace{4\baselineskip}\subsection{#1}}
\newcommand{\field}[1]{\par\textbf{#1}\par\noindent\rule{\linewidth}{0.3pt}\par\vspace{0.5em}}
\begin{document}
% Distinct page labels avoid duplicate Hyperref destinations after the cover.
\pagenumbering{roman}
\begin{titlepage}\thispagestyle{empty}\raggedright
\vspace*{0.65in}
{\sffamily\bfseries\color{teal} ARCHITECTURE DOCUMENTATION / CONFORMANCE GUIDE\par}
\vspace{0.35in}
{\sffamily\bfseries\color{navy}\Huge Architecture Description Conformance Guide\par}
\vspace{0.2in}
{\sffamily\Large\color{muted} Claim boundaries, clause-indexed evidence, assessment workflow and reporting\par}
\vspace{0.35in}\noindent{\color{teal}\rule{\linewidth}{1.5pt}}\par
\vspace{0.3in}
An edition-specific assessment guide grounded in the supplied 2022 standard. It separates AD, ADF, ADL, viewpoint and model-kind claims and avoids treating a template as certification.

\vfill
{\sffamily\small Edition 1.0 \quad / \quad 25 November 2026\par
Status: library working edition \quad / \quad Original engineering guidance\par}
\end{titlepage}
\pagenumbering{arabic}
\tableofcontents\clearpage
\docsection{Claim boundaries and assessment policy}
ISO 42010 clause 4 distinguishes five claim objects. The AD is assessed against clause 6, ADF against 7.1, ADL against 7.2, viewpoint against 8.1 and model kind against 8.2. Clause 8.3 governs definition of view methods when guidance is needed. Conformance is not a claim about architecture quality, implementation correctness or independent accreditation. Do not tailor away requirements. This guide uses original audit prompts and clause locators, not a replacement for normative wording.

\begin{longtable}{@{}>{\RaggedRight\arraybackslash}p{0.297\linewidth}>{\RaggedRight\arraybackslash}p{0.297\linewidth}>{\RaggedRight\arraybackslash}p{0.297\linewidth}@{}}
\toprule
\textbf{Claim object} & \textbf{Applicable clause} & \textbf{Evidence scope}\\\midrule
\endhead
AD & 6 & Versioned project description and referenced viewpoint specifications.\\
ADF & 7.1 & Framework specification, domain, conventions, coverage and adherence provisions.\\
ADL & 7.2 & Language specification, syntax/semantics, implemented model kinds and associated methods.\\
Viewpoint & 8.1 & Versioned convention specification, coverage and interpretation guidance.\\
Model kind & 8.2 & Model syntax and interpretation, methods, version and sources.\\
\bottomrule
\end{longtable}
\docsection{Assessment procedure}
\begin{itemize}
\item Freeze the exact claim object, version, scope, referenced artifacts and evidence commit. Obtain the licensed normative text for the assessor.
\item Distinguish mandatory provisions, recommendations and permissions. Expand each clause locator into individual normative tests in a controlled assessment workspace. Do not treat the grouped rows below as exhaustive atomic requirements.
\item For each test record result, evidence locator and reviewer. Mark a conditional provision inapplicable only with an explicit rationale grounded in the provision, never as discretionary tailoring.
\item Assess all unresolved violations and inconsistencies; distinguish a requirement to record an issue from a substantive unmet requirement. A recorded nonconformance does not convert failure into a pass.
\item Review any referenced viewpoint and model-kind specification. Do not infer that an AD claim automatically proves an ADF or ADL claim.
\item Issue a bounded conclusion only after every mandatory applicable test is supported. Preserve review history and reopen when the baseline changes.
\end{itemize}
\begin{longtable}{@{}>{\RaggedRight\arraybackslash}p{0.457\linewidth}>{\RaggedRight\arraybackslash}p{0.457\linewidth}@{}}
\toprule
\textbf{Status} & \textbf{Meaning}\\\midrule
\endhead
Pass & Evidence demonstrates the specific applicable obligation.\\
Fail & Evidence establishes unmet obligation.\\
Not verified & Evidence missing or assessment incomplete; blocks positive claim.\\
Conditional not applicable & Provision condition not met; record justification and review.\\
Recommendation & Separate advisory evaluation; not counted as mandatory failure by itself.\\
\bottomrule
\end{longtable}
\docsection{AD assessment evidence map}
\begin{longtable}{@{}>{\RaggedRight\arraybackslash}p{0.13\linewidth}>{\RaggedRight\arraybackslash}p{0.28\linewidth}>{\RaggedRight\arraybackslash}p{0.48\linewidth}@{}}
\toprule
\textbf{Clause} & \textbf{Assessment area} & \textbf{Evidence to inspect}\\\midrule
\endhead
6.1 & Identification and purpose & AD metadata, EoI/environment, purpose, project-specific information.\\
6.2 & Stakeholders and limitations & Stakeholder register; prior evaluation consideration; limitations statement; identified and explained nonconformances.\\
6.3 & Perspectives & Perspective holders, resulting concerns and identified perspectives.\\
6.4 & Concerns & Relevant concerns and associations to holders.\\
6.5 & Aspects & Relevant aspects and concern associations.\\
6.6 & Viewpoints & Versioned specifications or references; Clause 8 assessment; all concerns framed; perspectives mapped.\\
6.7 & Views & At least one view per used viewpoint; version and identity; coverage of perspectives, concerns and aspects; governing conventions; components and known issues.\\
6.8 & View components & Component version, viewpoint conformity, model kind if applicable or legend for non-model information.\\
6.9.1 & Consistency & Known inconsistencies recorded; consistency analysis considered.\\
6.9.2 & Correspondences & List; participants, AD identities and governing method(s).\\
6.9.3 & Methods & List and references; each applied method outcome or known violations.\\
6.10.1 & Decisions & Essential decisions recorded; decision strategy and alternatives considered.\\
6.10.2 & Rationale & Essential decision rationale; selection and limitations rationale considered.\\
\bottomrule
\end{longtable}
For clause 6.2, specifically review previous evaluations, resource limitations and explanations of nonconformance. For 6.7, verify every view addresses the concerns framed by its governing viewpoint, with scope clearly bounded; also verify global coverage of all identified concerns, perspectives and aspects. For 6.9.3, absence of an associated correspondence cannot be treated as proof a method holds.

\docsection{Viewpoint and model-kind assessment}
\begin{longtable}{@{}>{\RaggedRight\arraybackslash}p{0.457\linewidth}>{\RaggedRight\arraybackslash}p{0.457\linewidth}@{}}
\toprule
\textbf{Clause} & \textbf{Evidence prompts}\\\midrule
\endhead
8.1 & Associated perspectives, framed concerns, related aspects, typical stakeholders, model kinds and legends, within-view correspondence methods, sources; legend interpretation guidance; view methods when needed.\\
8.2 & Modeling conventions, associated view and correspondence methods, project/organization version identification and sources.\\
8.3 & When needed, specify construction/use methods; document transformations and representation of information sources, including non-model sources.\\
\bottomrule
\end{longtable}
A selected architectural style alone is insufficient as a viewpoint specification. A UML diagram alone does not establish the model kind conventions applicable to a given AD. Check that the specification and actual representation agree.

\docsection{ADF and ADL assessment boundaries}
If the organization elects to claim framework or language conformance, use a separate claim record and evaluate the complete applicable clause. Clause 7.1 distinguishes framework identification and content, recommended provisions and adherence rules. Review typical stakeholders, concerns and viewpoints; perspectives, aspects, structuring and legends as applicable; correspondence methods and rules for framework adherence. A local viewpoint catalog is not automatically an ADF. Clause 7.2 covers language concerns/aspects, view methods, model kinds, viewpoints if implemented, correspondence methods and version. A diagram renderer is not automatically an ADL.

\docsection{Evidence record and example conclusion}
\begin{longtable}{@{}>{\RaggedRight\arraybackslash}p{0.457\linewidth}>{\RaggedRight\arraybackslash}p{0.457\linewidth}@{}}
\toprule
\textbf{Field} & \textbf{Illustrative entry}\\\midrule
\endhead
Claim & AD-REF v1.0 / example only\\
Clause test & 6.6 / concern framing coverage\\
Evidence & Concern-to-viewpoint register in reference AD\\
Result & Pass for illustrative mapping check only; no full normative assessment\\
Reviewer & Unassigned\\
Finding & Runtime performance and security targets not tested\\
Disposition & No formal conformance claim\\
\bottomrule
\end{longtable}
A numerical percentage cannot substitute for a mandatory-requirement disposition. Report document scope, evidence reviewed, each unresolved failure, and whether conformance is demonstrated or not demonstrated. Keep technical quality findings and recommendations in separate categories. The companion review checklist covers usability and engineering value as well as evidence review.

\docsection{Source register and scope}
\begin{enumerate}
\item ISO/IEC/IEEE 42010:2022, \emph{Software, systems and enterprise --- Architecture description}, second edition, November 2022. Supplied licensed copy consulted; citations refer to clause numbers, not PDF pagination. Official record: \url{https://www.iso.org/standard/74393.html}.
\item P. Clements et al., \emph{Documenting Software Architectures: Views and Beyond}, second edition, Addison-Wesley, 2011, ISBN 978-0-321-55268-6. Supplied copy consulted; citations refer to book chapters and sections.
\end{enumerate}
All procedures, examples, record identifiers, tables, and proposed engineering controls in this suite are original implementation guidance. Neither source is redistributed. The book predates the 2022 edition; its historical standards comparison is not evidence of compliance with the later edition. Consult the actual standard for normative wording and applicability. This suite is not certification, legal advice, or a substitute for independent architecture evaluation.
\end{document}
